% 张三慧《大学物理学》第四版 V2，第24章《光的偏振》
% 主教材PDF物理页313--315（印刷296--298）；辅导书S PDF328--333。
% 图24.38--24.43存在原图几何依赖，专家图形复核待办。
\section*{第24章\quad 光的偏振}
\subsection*{章末习题24.1--24.16}
\paragraph{习题 24.1}
自然光通过两个偏振片，偏振方向夹角60°，
透过光强为 \(I_1\)。
在两片之间插入一片方向与前后两片各夹30°的
偏振片，求新的透射强度（不计吸收）。
\paragraph{习题 24.2}
自然光通过两片相互重叠的偏振片。
如果透射强度等于(1)其\emph{最大}透射强度的
三分之一，或(2)原入射自然光强度的三分之一，
分别求两片透振方向夹角。
\paragraph{习题 24.3}
两片偏振片 \(P_1,P_2\) 初始平行；
强度为 \(I_0\) 的自然光正入射第一片，
第二片先绕入射光轴转 \(\theta\)，
再绕竖直轴转 \(\varphi\)（原图24.38）。
求第二片透射光强，说明为何需用在横截面上的
投影角而非直接把 \(\theta,\varphi\) 的余弦相乘。
\paragraph{习题 24.4}
对原图24.39中的五种(a)--(e)情况，
已知以非偏振光或线偏振光入射两介质界面，
图中 \(i_b\) 为布儒斯特角且 \(i\ne i_b\)。
按菲涅耳定律画出各自反射/折射光的传播方向与
偏振状态（点与短线标记），对于起偏角
说明哪一种偏振分量没有反射光。
\paragraph{习题 24.5}
水折射率 \(n_w=1.33\)、玻璃 \(n_g=1.50\)。
求水到玻璃、玻璃到水反射时的起偏振角，
并说明两角的关系。
\paragraph{习题 24.6}
光在两种介质界面上的\emph{全反射临界角}
为45°，求从同一高折射率一侧入射时的起偏振角。
\paragraph{习题 24.7}
某透明介质在空气中的布儒斯特起偏振角
为 \(58^\circ\)，求其折射率。
\paragraph{习题 24.8}
从平静湖面反射的太阳光为完全线偏振光时，
太阳在地平线上方的仰角是多少？
取水的折射率1.33，空气近似为1。
\paragraph{习题 24.9}
方解石切割为正三角形棱镜，
光轴垂直于正三角形截面（图24.40）。
自然光从一面以入射角 \(i\) 进入，
要求 \(e\) 光在棱镜内部沿平行底面的方向传播。
求 \(i\)，并画出 \(o\) 光传播方向及两种光振动方向。
可用原书给出的 \(n_o\simeq1.6584,n_e\simeq1.4864\)。
\paragraph{习题 24.10}
原图24.41：由两块各带45°角的方解石棱镜
\(ABD\) 和 \(BCD\) 合成方形棱镜，
第一棱镜光轴与 \(AB\) 平行，
第二棱镜光轴垂直图面。
自然光垂直入射 \(AB\)。
按原图画 \(o,e\) 光在两棱镜中传播方向，
及两束出射光的线偏振振动方向。
\paragraph{习题 24.11}
原图24.42：\(P_1,P_2\) 为两片\emph{正交}线偏振片，
中间有四分之一波片 \(C\)，其光轴与
\(P_1\) 透振方向夹60°。
强度 \(I_i\) 的单色自然光垂直入射 \(P_1\)。
求图示①（经 \(P_1\)）、②（经 \(C\)）、
③（经 \(P_2\)）各区的偏振态和光强。
\paragraph{习题 24.12}
晶体对波长632.8nm的光的主折射率
\(n_o=1.66,n_e=1.49\)。
制作该波长的四分之一波片，求最小正厚度，
并说明晶体光轴相对波片表面应如何取向。
\paragraph{习题 24.13}
某石英四分之一波片为波长800nm设计，
假设 \(n_o,n_e\) 与波长无关。
波长400nm的线偏振光入射，
振动方向与波片光轴夹45°。
求透射光偏振态及主轴方向。
\paragraph{习题 24.14}
原图24.43的尼科耳棱镜由两块直角方解石棱镜
以折射率1.55的加拿大胶接合，
棱镜底角68°，底面涂黑。
自然光从平行底面的方向入射，证明
\(o\) 光在胶层全反射并被吸收，
\(e\) 光透过胶层和第二棱镜，求透射光偏振态。
参考折射率 \(n_o=1.658,n_e\simeq1.486\)。
\paragraph{习题 24.15}
石英对于波长396.8nm的右旋圆偏振光
折射率 \(n_R=1.55810\)，对于左旋圆偏振光
\(n_L=1.55821\)。
求该波长的旋光率（单位长度旋转角，注明
观看传播方向的正负号约定）。
\paragraph{习题 24.16}
两列等强、等频激光分别沿
\(+x\) 与 \(-x\) 传播，
其电振动沿 \(y\) 与 \(z\) 方向，
在 \(x=0\) 合成为与 \(y\) 轴夹45°的线偏振光。
求沿 \(+x\) 方向每隔 \(\lambda/8\)
的合成偏振态，区分线偏振、左右旋圆偏振；
用相位表达式说明周期性。
